\subsection{与复合算子相联系的阿佩尔多项式}

给定一个复合算子 \(U = \sum_{k=0}^{\infty} \alpha_k \mathrm{D}^k \neq 0\) ，设 \(p\) 为整数 \(k\) （使 \(\alpha_k \neq 0\) 的）中最小的一个；我们称 \(p\) 为算子 \(U\) 的阶。

命题 2. 每个 0 阶复合算子在复合算子代数 \(\Gamma\) （K{[}X{]} 上的）中可逆。

事实上，形式级数 \(\sum_{k=0}^{\infty}\alpha_{k}S^{k}\) （满足 \(\alpha_{0}\neq0\) ）在环 K{[}{[}S{]}{]} 中可逆（Alg., IV. 41）；于是命题由 VI, p.~271 的定理 1 得出。

命题 3. 设 U 是一个 p 阶复合算子；则 \(U(f) = 0\) （对每个次数 \textless{} p 的多项式 f）；对每个多项式 \(f \neq 0\) （次数 \(n \geqslant p\) 的），\(U(f)\) 是一个次数为 n - p 的多项式 \(\neq 0\) 。

这是 VI, p.~271 的公式 (2) 与 \(U\) 的阶的定义的直接推论。

显然，每个算子 \(U\) （\(p\) 阶的）都可以以唯一的方式写成 \(U = \mathrm{D}^p V = V\mathrm{D}^p\) ，其中 \(V\) 是一个 0 阶算子（从而可逆）。

定义 2. 设 \(U = \mathrm{D}^p V\) 是一个 \(p\) 阶复合算子（\(\mathrm{K}[\mathrm{X}]\) 上的）。多项式 \(u_n(\mathrm{X}) = V^{-1}(\mathrm{X}^n)\) 称为指标为 \(n\) 而与算子 \(U\) 相联系的阿佩尔多项式。

\[
\text {   If   } V ^ {- 1} = \sum_ {k = 0} ^ {\infty} \frac {1}{k !} \beta_ {k} D ^ {k} (\text {   with   } \beta_ {0} \neq 0) \text {   one   thus   has   }
\]

\[
u _ {n} (\mathrm{X}) = \sum_ {k = 0} ^ {n} {\binom {n} {k}} \beta_ {k}   \mathrm{X} ^ {n - k}.\tag{6}
\]

可以验证 \(u_{n}\) 是一个 \(n\) 次多项式（prop. 3）；此外

\[
u _ {n} (0) = \beta_ {n}.
\]

命题 4. 与 \(U\) 相联系的阿佩尔多项式满足下列关系

\[
\frac {d u _ {n}}{d \mathrm{X}} = n u _ {n - 1}\tag{7}
\]

\[
u _ {n} (\mathrm{X} + \mathrm{Y}) = \sum_ {k = 0} ^ {n} {\binom {n} {k}} u _ {n - k} (\mathrm{X})   \mathrm{Y} ^ {k}\tag{8}
\]

\[
U \left(u _ {n} (\mathrm{X})\right) = \frac {n !}{(n - p) !} \mathrm{X} ^ {n - p}.\tag{9}
\]

这些公式事实上分别等价于下列关系（记住定义 2）：

\[
\mathbf {D} V ^ {- 1} = V ^ {- 1} \mathbf {D}\tag{10}
\]

\[
\left(\exp \left(\mathrm{YD} _ {\mathrm{X}}\right)\right) V _ {\mathrm{X}} ^ {- 1} = V _ {\mathrm{X}} ^ {- 1} \exp \left(\mathrm{YD} _ {\mathrm{X}}\right)\tag{11}
\]

\[
U V ^ {- 1} = \mathbf {D} ^ {p}.\tag{12}
\]

命题 5. 对每个多项式 \(f \in \mathrm{K}[\mathrm{X}]\) 和每个复合算子 \(U\) （\(p\) 阶的），有

\[
f ^ {(p)} (\mathrm{X} + \mathrm{Y}) = \sum_ {k = 0} ^ {\infty} \frac {1}{k !} U (f ^ {(k)} (\mathrm{X})) u _ {k} (\mathrm{Y})\tag{13}
\]

（广义泰勒展开）。

事实上，若令 \(U = \mathrm{D}^p V = V\mathrm{D}^p\) ，则有（VI, p.~270, 公式 (1)）

\[
V _ {\mathrm{X}} ^ {- 1} (f (\mathrm{X} + \mathrm{Y})) = V _ {\mathrm{Y}} ^ {- 1} (f (\mathrm{X} + \mathrm{Y})) = \sum_ {k = 0} ^ {\infty} \frac {1}{k !} f ^ {(k)} (\mathrm{X}) u _ {k} (\mathrm{Y})\tag{14}
\]

这是由泰勒公式及 VI, p.~273 的定义 2 得到的；只需对公式 (14) 的首末两项应用算子 \(U_{\mathrm{X}}\) 即得 (13)。

